\documentclass{article}

\usepackage{float}
\usepackage{longtable}
\restylefloat{table}

\usepackage{booktabs}

\title{Team Productivity: POC\\\progname}

\author{\authname}

\date{}

\input{../Comments.text}
\input{../Common.text}

\begin{document}

\maketitle

This document summarizes the contributions of each team member up to the POC
Demo.  The time period of interest is the time between the beginning of the term
and the POC demo.

\wss{Please do not delete sections or modify the table formats.  Standardization
helps the instructors and TAs to review these documents. You may add sections to
the report.}

\wss{Remove the instructor guidelines/comments in the final version.}

\wss{Fill in the requested information, like project url, TA name, supervisor name, etc.}

\section{Demo Plans}

\wss{What is your feasibility risk?}

\wss{What will you be demonstrating}

\section{Team Meeting Attendance}

\begin{table}[H]
\centering
\begin{tabular}{ll}
\toprule
\textbf{Student} & \textbf{Meetings}\\
\midrule
Total & 8\\
Hemal Rajput & 8\\
Gurmanjot Minhas & 8\\
James Araiche & 8\\
Lily Porter & 8\\
Sarah Dorfman & 8\\
\bottomrule
\end{tabular}
\end{table}

One extra team meeting was conducted between James and Sarah to discuss research for the machine learning portion of the project. This task was explicitly assigned to them so no other team members needed to be in attendance. Therefore, it is excluded from the total count to clearly indicate the attendance for meetings held for the whole team.

\section{Supervisor/Stakeholder Meeting Attendance}

\noindent \textbf{Supervisor's Name: } Dr. Lyn Turkstra, Dr. Daniel Pape

\begin{table}[H]
\centering
\begin{tabular}{ll}
\toprule
\textbf{Student} & \textbf{Meetings}\\
\midrule
Total & 4\\
Hemal Rajput & 4\\
Gurmanjot Minhas & 4\\
James Araiche & 4\\
Lily Porter & 4\\
Sarah Dorfman & 4\\
\bottomrule
\end{tabular}
\end{table}

\section{Lecture Attendance}

\begin{table}[H]
\centering
\begin{tabular}{ll}
\toprule
\textbf{Student} & \textbf{Lectures}\\
\midrule
Total & 12\\
Hemal Rajput & 11\\
Gurmanjot Minhas & 11\\
James Araiche & 11\\
Lily Porter & 11\\
Sarah Dorfman & 11\\
\bottomrule
\end{tabular}
\end{table}

Majority of us missed the VnV plan lecture due the fact it was midterm season. Sarah had attended the VnV lecture and was able to fill us in on important information that we missed.

\section{TA Document Discussion Attendance}

\wss{For each team member how many of the informal document discussion meetings
with the TA were attended over the time period of interest. The first row is the
total number of discussion meetings scheduled by the TA with the team.}

\noindent \textbf{TA's Name: } [fill in this information]

\begin{table}[H]
\centering
\begin{tabular}{ll}
\toprule
\textbf{Student} & \textbf{Meetings}\\
\midrule
Total & Num\\
Name 1 & Num\\
Name 2 & Num\\
Name 3 & Num\\
Name 4 & Num\\
Name 5 & Num\\
\bottomrule
\end{tabular}
\end{table}

\wss{If needed, an explanation for the attendance can be provided here.}

\section{Commits, Lines Added and Lines Deleted, Pull Requests}

\wss{For each team member how many commits to the main branch have been made
over the time period of interest.  The total is the total number of commits for
the entire team since the beginning of the term.  The percentage is the
percentage of the total commits made by each team member.}

\wss{There are also columns for lines added and lines deleted to give a sense of
the magnitude of each team member's contributions.  Note that these numbers can
be somewhat misleading, as a team member who does refactoring may have high
numbers in both columns, while a team member who implements a large new feature
may have a high number of lines added, but few lines deleted.  As for all commit
information, these numbers will be interpreted with care.}

\wss{The pull requests column counts the number of pull requests created by each
team member.  This is a measure of how many contributions were made in branches,
which may not be reflected in the commit counts for the main branch.}

\begin{table}[H]
\centering
\begin{tabular}{lllllll}
\toprule
\textbf{Student} & \textbf{Commits} & \textbf{Percent} & \textbf{Lines Added} &
\textbf{Lines Deleted} & \textbf{Pull Requests}\\
\midrule
Total & Num & 100\% & total added & total deleted & total PRs\\
Name 1 & Num & \% & num & num & num\\
Name 2 & Num & \% & num & num & num\\
Name 3 & Num & \% & num & num & num\\
Name 4 & Num & \% & num & num & num\\
Name 5 & Num & \% & num & num & num\\
\bottomrule
\end{tabular}
\end{table}

\wss{If needed, an explanation for the counts can be provided here.  For
instance, if a team member has more commits to unmerged branches, these numbers
can be provided here.  If multiple people contribute to a commit, git allows for
multi-author commits.}

\section{Issue Tracker}

\wss{For each team member how many issues have they authored (including open and
closed issues (O+C)) and how many have they been assigned (only counting closed
issues (C only)) over the time period of interest.}

\begin{table}[H]
\centering
\begin{tabular}{lll}
\toprule
\textbf{Student} & \textbf{Authored (O+C)} & \textbf{Assigned (C only)}\\
\midrule
Name 1 & Num & Num \\
Name 2 & Num & Num \\
Name 3 & Num & Num \\
Name 4 & Num & Num \\
Name 5 & Num & Num \\
\bottomrule
\end{tabular}
\end{table}

\wss{If needed, an explanation for the counts can be provided here.}

\section{CICD via GitHub Actions}

\wss{Say how CICD technology is used and will be used in your project.}

\wss{Provide links to your CICD yaml files}

\section{Custom Documents (CDs)}

As documented in TeamComposition.csv, our original plan was to submit a User/Stakeholder/Client interview report as the fall course deliverable (CD) and a usability report in the winter term. Since then, we have learned that our supervisors prefer to review mockups of PronouncedFish, as this lets them give feedback and shape the app to fit their clinical workflow. We have drafted and iterated on several mockups and presented them to our supervisors for comment. Because of this, we have replaced the interview report with a wireframe report as our fall CD. The report captures the main features of the app and the supervisor feedback on its layout and functionality. For example, the mouth-placement map and progress-over-time graphs in the session analysis wireframes were added after supervisors said these visualizations would be useful. Our plan for the winter CD is unchanged: a usability report. By then we expect to have a working prototype in the hands of our supervisors and their students, which will give us valuable feedback for improving the project.

\section{Team Charter Trigger Items}

The following triggers are taken from the team charter.

\begin{itemize}
    \item[\textbf{Attendance:}] Missing a team meeting without notifying
    the team in advance (1--3 days for personal or school conflicts) is a
    trigger. Attendance is tracked at every meeting, with the following
    thresholds:
    \begin{itemize}
        \item Below 80\%: the member is flagged after each meeting.
        \item Below 60\%: a team discussion is held and an improvement
        plan may be created.
        \item Below 50\%: a meeting with the TA is arranged.
        \item Below 30\%: a meeting with Dr.\ Smith is arranged.
    \end{itemize}
    Arriving without required meeting materials, without notifying the
    team beforehand, is also a trigger.

    \item[\textbf{Quality:}] A pattern of failing to deliver quality work
    for meetings or deliverables is a trigger. The team discusses next
    steps at a team meeting and, if the pattern continues, raises it with
    Dr.\ Smith.

    \item[\textbf{Attitude:}] A violation of the code of conduct is a
    trigger. The issue is first discussed between the individuals
    involved, then at the team level, then with the TAs or professor. Serious
    or harmful conflicts go directly to Dr.\ Smith.

    \item[\textbf{Stay on track:}] Missing two project deadlines triggers
    frequent or daily check-ins. If deadlines continue to be missed, the
    team will arrange a discussion with Dr.\ Smith.

    \item[\textbf{Workload:}] An uneven distribution of GitHub issues
    among team members is a trigger. The team will reconvene to
    redistribute tasks.
\end{itemize}

\begin{longtable}[c]{| c | c |}
\caption{Team Charter Violations.\label{CharterViolations}}\\
    \hline
     \multicolumn{2}{| c |}{No violations to report!}\\
     \hline
     \endfirsthead
 \end{longtable}

Since the team has no violations to report, there does not need to be any changes made to the team charter.
 
\section{Additional Productivity Metrics}

\wss{If your team has additional metrics of productivity, please feel free to
add them to this report.  If not, please explicitly state that there are no additional
metrics.}

\wss{Additional metrics can include things like code reviews done, pull requests
created, count of joining meetings late, count of number of times contributions
had to be corrected, number of internal deadlines missed, test cases written, etc.}

\wss{We are looking for data on these metrics, not just a list of additional
metrics the team is planning on tracking.  However, if all you have is a plan,
please share it here.}

\end{document}